\documentclass[a4paper,12pt]{article}
\usepackage[english]{babel}
\usepackage{hyperref}
\hypersetup{hidelinks}
\usepackage{amsmath}
\usepackage{amssymb}
\usepackage[varvw]{newtxmath}

% page geometry — identical to thesis
\setlength{\oddsidemargin}{8mm}
\setlength{\textwidth}{145mm}
\setlength{\topmargin}{2mm}
\setlength{\textheight}{225mm}
\setlength{\headheight}{0mm}
\setlength{\parskip}{0.5mm}
\setlength{\parindent}{9mm}
\linespread{1.12}
\renewcommand{\floatpagefraction}{.93}
\renewcommand{\textfraction}{0.02}

% ─────────────────────────────────────────────────────────────────────────────
\title{{\bf \Large TDSE-CN — Development Log}}
\author{Ludovico Fecia di Cossato}
\date{}

\begin{document}

\maketitle
\thispagestyle{empty}

\noindent\rule{\textwidth}{0.4pt}
\vskip 2mm
\noindent\textit{Session-by-session record of design decisions and implementation
progress. Kept separate from the main article so that the official documentation
remains clean and version-independent.}
\vskip 2mm
\noindent\rule{\textwidth}{0.4pt}

% ─────────────────────────────────────────────────────────────────────────────
\section*{21 March 2026 — Session 1}
\textbf{Theory, project structure, input file, params module.}

\begin{itemize}

  \item Defined the physical scenario: 1D Gaussian wave packet propagating
        toward a rectangular potential barrier, studied via the TDSE with
        the Crank--Nicolson scheme.

  \item Established the full theoretical roadmap in five blocks (0--5),
        from the TDSE to the computation of observables
        (norm, $\langle x\rangle$, $T$, $R$).

  \item Discussed space-time discretisation: spatial grid
        $x_j = j\,\Delta x$ and temporal grid $t_n = n\,\Delta t$;
        $\psi(x,t)$ becomes a complex vector at each time step.
        Introduced the centred finite-difference approximation for the
        second spatial derivative and identified the time-level choice as
        the key design decision (leading to Crank--Nicolson in block~2).

  \item Defined the project directory structure:
        \texttt{input/}, \texttt{src/}, \texttt{output/} (gitignored),
        \texttt{analysis/}, \texttt{docs/}, \texttt{makefile}.

  \item Wrote \texttt{input/input.txt} with all physical and numerical
        parameters in atomic units
        (\texttt{L}, \texttt{dx}, \texttt{dt}, \texttt{T\_tot},
         \texttt{x0}, \texttt{sigma}, \texttt{k0},
         \texttt{x\_bar}, \texttt{w\_bar}, \texttt{V0}).
        Added inline unit annotations for each parameter.
        Decision: all source comments in English.

  \item Implemented \texttt{src/params.hpp}: \texttt{Params} struct with
        fields for all input parameters plus derived quantities $N$
        (grid points) and $M$ (time steps); function declarations using
        \texttt{\#pragma once} (chosen over \texttt{\#ifndef} guards for
        modern C++).

  \item Implemented \texttt{src/params.cpp}: \texttt{read\_params()} reads
        key--value pairs line by line, skipping \texttt{\#} comments;
        computes $N$ and $M$ via \texttt{std::round} to avoid
        floating-point truncation. \texttt{print\_params()} formats a
        labelled summary to stdout.

  \item Written minimal \texttt{src/tdse.cpp} calling
        \texttt{read\_params()} and \texttt{print\_params()} to test the
        I/O pipeline. Compiled successfully with \texttt{makefile} using
        \texttt{g++ -std=c++17 -Wall -Wextra -O2}.

  \item Produced project documentation:
        \texttt{TDSE\_CN\_article.tex} (official article, no updates
        section) and \texttt{TDSE\_CN\_updates.tex} (this file).

\end{itemize}

% ─────────────────────────────────────────────────────────────────────────────
% Template for future entries — copy and fill in:
%
% \section*{DD Month YYYY — Session N}
% \textbf{Short description of the session.}
%
% \begin{itemize}
%   \item ...
% \end{itemize}
%
% ─────────────────────────────────────────────────────────────────────────────

\end{document}